\documentclass[10pt,a4paper]{amsart}
\usepackage[margin=19mm]{geometry}
\usepackage{amsmath,amssymb,mathtools}
\usepackage[scr=rsfs,cal=euler]{mathalfa}
\usepackage{xfrac}
\usepackage[unicode,hidelinks,bookmarks=false]{hyperref}
\newcommand{\ie}{i.e.\ }
\newcommand{\cf}{cf.\ }
\newcommand{\resp}{respectively\ }
\newcommand{\Subsection}{\subsection}
\providecommand{\nobreakdash}{\nobreak\hbox{-}}
\setlength{\emergencystretch}{2em}
\multlinegap0pt
\newtheorem{lemma}{Lemma}
\newtheorem{proposition}{Proposition}
\newtheorem{theorem}{Theorem}
\newtheorem{corollary}{Corollary}
\usepackage{biblatex}
\begin{document}
\title[Computing lower expectations with respect to total variation]{Computing lower expectations with respect to total variation distance and chi-squared divergence balls}
\author{Jasper De Bock}
\date{}
\maketitle
\noindent\emph{Source and licence.} Computing lower expectations with respect to total variation distance and chi-squared divergence balls. Version arXiv:2605.30091v1 (CC BY 4.0). This material is distributed under the Creative Commons Attribution 4.0 International licence, http://creativecommons.org/licenses/by/4.0/, as recorded in the arXiv OAI licence metadata for this exact version and shown on the abstract page https://arxiv.org/abs/2605.30091v1. Full rights notice: licenses/arxiv-2605.30091-CC-BY-4.0.txt.\par\smallskip
\noindent\textit{Editorial scope.} Counted unit: the original \texttt{theorem} environment \texttt{thm:lowerexpectationtv} at source lines 58--83 together with its complete proof at lines 159--220; the delivered document keeps source lines 43--221 so that every referenced label and every citation resolves inside it. Native editable source is the paper's own concatenated TeX; the AMS wrapper is editorial formatting only. No publisher class, upstream script or unrelated artwork is redistributed.
\section{Total variation balls}\label{sec:tv}

The first case that we consider is when $d$ is the total variation distance, which is defined for all $q\in\Sigma_\mathcal{X}$ by
\begin{equation*}
    d_\mathrm{TV}(q,p):=\frac{1}{2}\sum_{x\in\mathcal{X}}|q(x)-p(x)|.
\end{equation*}
In this case, for any $\delta\in[0,1]$, we write $B_\mathrm{TV}(p,\delta)$ instead of $B_{d_\mathrm{TV}}(p,\delta)$.

For any given $f\in\mathbb{R}^\mathcal{X}$, our aim is to compute $\underline{E}_{B_\mathrm{TV}(p,\delta)}(f)$, which we will henceforth denote by $\underline{E}^\mathrm{TV}_{p,\delta}(f)$ for the sake of brevity.

The first step in computing $\underline{E}^\mathrm{TV}_{p,\delta}(f)$ is to sort the elements of $\mathcal{X}=\{x_1,\ldots,x_n\}$ in such a way that $f(x_1)\leq f(x_2)\leq\cdots\leq f(x_n)$. We will henceforth assume that $\mathcal{X}$ is ordered in this way.

Now let $r\in\{1,\ldots,n\}$ be the smallest index such that $\delta\geq\sum_{i=r+1}^{n} p(x_i)$. Note that such an index must exist, since $\sum_{i=n+1}^n p(x_i)=0$ and $\delta\in[0,1]$.
The following result then provides the desired closed form expression.


\begin{theorem}\label{thm:lowerexpectationtv}
Let $r$ be defined as above. Then
\begin{equation*}
    \underline{E}^\mathrm{TV}_{p,\delta}(f)
    =
    E_{q_\mathrm{TV}}(f)
\end{equation*}
where, if $r>1$, $q_\mathrm{TV}\in\Sigma_\mathcal{X}$ is given by
\begin{equation*}
    q_\mathrm{TV}(x_i):=
    \begin{cases}
        p(x_1)+\delta & \text{if } i=1,\\
        p(x_i) & \text{if } i\in\{2,\ldots,r-1\},\\
        \sum_{i=r}^{n} p(x_i)-\delta & \text{if } i=r,\\
        0 & \text{if } i\in\{r+1,\ldots,n\}.
    \end{cases}
\end{equation*}
and, if $r=1$, $q_\mathrm{TV}\in\Sigma_\mathcal{X}$ is given by
\begin{equation*}
    q_\mathrm{TV}(x_i):=
    \begin{cases}
        1 & \text{if } i=1,\\
        0 & \text{if } i\in\{2,\ldots,n\}.
    \end{cases}
\end{equation*}
\end{theorem}

Intuitively, this corresponds to constructing $q_\mathrm{TV}$ starting from $p$, by moving as much mass as possible from the elements of $\mathcal{X}$ with the highest $f$-values to the element of $\mathcal{X}$ with the lowest $f$-value, until we have moved a total mass of $\delta$ or, if $r=1$, have moved all available mass.

\section{chi-squared divergence balls}\label{sec:chi2}

The second case that we consider is when $d$ is the chi-squared divergence, often denoted as $\chi^2$-divergence, which, if $p(x_i)>0$ for all $i\in\{1,\ldots,n\}$, is defined for all $q\in\mathbb{R}^\mathcal{X}$ by
\begin{equation*}
    d_{\chi^2}(q,p):=\sum_{x\in\mathcal{X}}\frac{(q(x)-p(x))^2}{p(x)}
    =\sum_{x\in\mathcal{X}}p(x)\left(\frac{q(x)}{p(x)}-1\right)^2.
\end{equation*}
In this case, for any $\delta\geq0$, we write $B_{\chi^2}(p,\delta)$ instead of $B_{d_{\chi^2}}(p,\delta)$.

For any given $f\in\mathbb{R}^\mathcal{X}$, our aim is to compute $\smash{\underline{E}_{B_{\chi^2}(p,\delta)}(f)}$, which we will henceforth denote by $\underline{E}^{\chi^2}_{p,\delta}(f)$ for the sake of brevity.

The first step is again to sort the elements of $\mathcal{X}=\{x_1,\ldots,x_n\}$ in such a way that $f(x_1)\leq f(x_2)\leq\cdots\leq f(x_n)$. Let $\ell$ be the largest index such that $f(x_\ell)=f(x_1)$. Next, for any $k\in\{\ell,\ldots,n\}$, we define
\begin{equation*}
    m_k:=\sum_{i=1}^k p(x_i),~~
\mu_k:=\frac{1}{m_k}\sum_{i=1}^k p(x_i)f(x_i)
\text{ and }\sigma_k^2:=\frac{1}{m_k}\sum_{i=1}^k p(x_i)(f(x_i)-\mu_k)^2
\end{equation*}
We also let $\delta_\ell:=+\infty$ and, for any $k\in\{\ell+1,\ldots,n\}$, let
\begin{equation*}
    \delta_k:=\frac{1}{m_k}\left(\frac{\sigma_k^2}{(f(x_k)-\mu_k)^2}+1-m_k\right).
\end{equation*}
These $\delta_k$ are the critical values of $\delta$ at which the lower expectation $\underline{E}^{\chi^2}_{p,\delta}(f)$ changes its form, as we will see below. We start out by showing that they are nicely ordered.
\begin{proposition}\label{prop:orderingdeltak}
The $\delta_k$ are well-defined, positive and non-increasing in $k$:
\begin{equation*}
    +\infty=\delta_\ell>\delta_{\ell+1}\geq\cdots\geq\delta_{n-1}\geq\delta_n>0.
\end{equation*}
\end{proposition}

Now let $r\in\{\ell,\ldots,n\}$ be the largest index such that $\delta_r>\delta$. Due to the previous proposition, this means that $r$ is the unique value in $\{\ell,\ldots,n\}$ such that
\begin{equation*}
\delta_\ell\geq\cdots\geq\delta_r>\delta\geq\delta_{r+1}\geq\cdots\geq\delta_n\geq0
\end{equation*}

The following theorem then gives us the desired closed form expression for the lower expectation that we are after.

\begin{theorem}\label{thm:lowerexpectationchi2}
With $r\in\{\ell,\ldots,n\}$ defined as above, we have that
\begin{equation}\label{eq:lowerexpectationchi2}
    \underline{E}^{\chi^2}_{p,\delta}(f)
    =
    \begin{cases}
    \mu_r-\sigma_r\sqrt{m_r\delta-(1-m_r)} & \text{if } r>\ell,\\
    f(x_1) & \text{if } r=\ell.
    \end{cases}
\end{equation}
\end{theorem}


If $\vert\mathcal{X}\vert=3$, Equation~\eqref{eq:lowerexpectationchi2} for example results in the following case distinction:
\begin{equation}\label{eq:lowerexpectationchi2three}
    \underline{E}^{\chi^2}_{p,\delta}(f)=
    \begin{cases}
        \mu_3-\sigma_3\sqrt{\delta} & \text{if } \delta<\delta_3,\\
        \mu_2-\sigma_2\sqrt{m_2\delta-(1-m_2)} & \text{if } \delta_3\leq\delta<\delta_2,\\
        f(x_1) & \text{if } \delta\geq\delta_2.
    \end{cases}
\end{equation}
If $\vert\mathcal{X}\vert=2$, the resulting case distinction can be simplified even further, resulting in the following simple expression:
\begin{equation}\label{eq:lowerexpectationchi2two}
    \underline{E}^{\chi^2}_{p,\delta}(f)=
    \begin{cases}
        \displaystyle\sum_{i=1}^2p(x_i)f(x_i)-\sqrt{\delta p(x_1)p(x_2)}\vert f(x_2)-f(x_1)\vert & \text{if } \delta<\frac{p(x_2)}{p(x_1)},\\
        f(x_1) & \text{if } \delta\geq\frac{p(x_2)}{p(x_1)}.
    \end{cases}
\end{equation}

\printbibliography

\appendix

\section{Proofs for the results in Section~\ref{sec:tv}}


\begin{proof}[Proof of Theorem~\ref{thm:lowerexpectationtv}]
If $r=1$, then $\delta\geq\sum_{i=2}^n p(x_i)$, which implies that
\begin{align*}
d_{\mathrm{TV}}(q_\mathrm{TV},p)
&=\frac{1}{2}\sum_{i=1}^n|q_\mathrm{TV}(x_i)-p(x_i)|\\
&=\frac{1}{2}\left(1-p(x_1)+\sum_{i=2}^n p(x_i)\right)
=\sum_{i=2}^n p(x_i)\leq\delta.
\end{align*}
The result then follows because $\min f\leq\underline{E}^\mathrm{TV}_{p,\delta}(f)\leq E_{q_\mathrm{TV}}(f)=f(x_1)=\min f$. So we can now assume that $r>1$, which implies that 
\begin{equation*}
\sum_{i=r+1}^n p(x_i)\leq\delta<\sum_{i=r}^{n} p(x_i)\leq\sum_{i=2}^n p(x_i)=1-p(x_1).
\end{equation*}

Since $B_\mathrm{TV}(p,\delta)$ is a non-empty closed subset of $\Sigma_\mathcal{X}$, the minimum $\underline{E}^\mathrm{TV}_{p,\delta}(f)$ is well-defined and there is at least one $q\in B_\mathrm{TV}(p,\delta)$ such that $\underline{E}^\mathrm{TV}_{p,\delta}(f)=E_q(f)$. Since the set of all $q$ that attain this minimum is also closed, we can choose $q^*$ to be one among them that maximizes $q(x_1)$.


We first show that $q^*(x_1)\geq p(x_1)$. Assume \emph{ex absurdo} that $q^*(x_1)<p(x_1)$. Since $q^*$ and $p$ are both probability mass functions, this implies that there is some $i\in\{2,\ldots,n\}$ such that $q^*(x_i)>p(x_i)$. Consider any such index. Then we can construct a new probability mass function $q'$ by mass $\alpha=\min\{p(x_1)-q^*(x_1),q^*(x_{i})-p(x_{i})\}>0$ from $x_{i}$ to $x_1$. This mass transfer does not alter the total variation distance to $p$, since we are moving the same mass $\alpha$ in opposite directions, and the expectation of $f$ does not increase since $f(x_1)\leq f(x_i)$. This contradicts the definition of $q^*$, which was chosen to have the highest value of $q(x_1)$ among all $q$ such that $\underline{E}^\mathrm{TV}_{p,\delta}(f)=E_q(f)$. So we must have that $q^*(x_1)\geq p(x_1)$, as claimed.

Next, we show that $q^*(x_i)\leq p(x_i)$ for all $i\in\{2,\ldots,n\}$. Assume \emph{ex absurdo} that there is some $i\in\{2,\ldots,n\}$ such that $q^*(x_i)>p(x_i)$. Consider any such index. Then we can construct a new probability mass function $q'$ by moving mass $\alpha=q^*(x_i)-p(x_i)>0$ from $x_i$ to $x_1$. This mass transfer does not alter the total variation distance to $p$, since we are moving the same mass $\alpha$ in opposite directions, and the expectation of $f$ does not increase since $f(x_1)\leq f(x_i)$. This contradicts the definition of $q^*$, which was chosen to have the highest value of $q(x_1)$ among all $q$ such that $\underline{E}^\mathrm{TV}_{p,\delta}(f)=E_q(f)$. So we must have that $q^*(x_i)\leq p(x_i)$ for all $i\in\{2,\ldots,n\}$, as claimed.

As a next step, we show that $q^*(x_1)=p(x_1)+\delta$. Assume \emph{ex absurdo} that $p(x_1)\leq q^*(x_1)<p(x_1)+\delta$. Since $\delta<1-p(x_1)$, this implies that $q^*(x_1)<1$. So there is some $i\in\{2,\ldots,n\}$ such that $0<q^*(x_i)\leq p(x_i)$. Consider any such $i$. Then we can construct a new probability mass function $q'$ by moving mass $\alpha=q^*(x_i)>0$ from $x_i$ to $x_1$. This mass transfer again does not alter the total variation distance to $p$, nor does it increase expectation of $f$, once more contradicting the definition of $q^*$. So we must have that $q^*(x_1)=p(x_1)+\delta$, as claimed.

We conclude that $q*$ belongs to the set
\begin{equation*}
\mathcal{Q}^*:=\{q\in\Sigma_\mathcal{X}: q(x_1)=p(x_1)+\delta\text{ and }q(x_i)\leq p(x_i)\text{ for all }i\in\{2,\ldots,n\}\}.
\end{equation*}
Since $d_\mathrm{TV}(q,p)=\delta$ for all $q\in\mathcal{Q}^*$ and therefore $\mathcal{Q}^*\subseteq B_\mathrm{TV}(p,\delta)$, and since $q^*$ was such that $\underline{E}^\mathrm{TV}_{p,\delta}(f)=E_{q^*}(f)$, this implies that $\underline{E}^\mathrm{TV}_{p,\delta}(f)=\min_{q\in\mathcal{Q}^*} E_q(f)$. Since the set of all $q\in\mathcal{Q}^*$ that attain the minimum $\min_{q\in\mathcal{Q}^*} E_q(f)$ is closed, we can choose $q^{**}$ to be one among them that maximizes $\sum_{i=2}^r q(x_i)$.

We now show that $q^{**}(x_i)=0$ for all $i\in\{r+1,\ldots,n\}$. Assume \emph{ex absurdo} that there is some $i\in\{r+1,\ldots,n\}$ such that $q^{**}(x_i)>0$. Consider any such index $i$. Then
\begin{align*}
\sum_{i=2}^r q^{**}(x_i)
&=1-q^{**}(x_1)-\sum_{i=r+1}^n q^{**}(x_i)\\
&<1-q^{**}(x_1)=1-p(x_1)-\delta\\
&\leq1-p(x_1)-\sum_{i=r+1}^n p(x_i)=\sum_{i=2}^r p(x_i).
\end{align*}
Since $q^{**}(x_i)\leq p(x_i)$ for all $i\in\{2,\ldots,n\}$, this implies that there is some $j\in\{2,\ldots,r\}$ such that $q^{**}(x_j)<p(x_j)$. Consider any such index $j$. Then we can construct a new probability mass function $q'$ by moving mass $\alpha=\min\{p(x_j)-q^{**}(x_j),q^{**}(x_i)\}>0$ from $x_i$ to $x_j$. This mass function $q'$ still belongs to $\mathcal{Q}^*$, and the expectation of $f$ does not increase since $f(x_j)\leq f(x_i)$. This however contradicts the definition of $q^{**}$, which was chosen to have the highest value of $\sum_{i=2}^r q(x_i)$ among all $q\in\mathcal{Q}^*$ that attain the minimum $\min_{q\in\mathcal{Q}^*} E_q(f)$. So we must have that $q^{**}(x_i)=0$ for all $i\in\{r+1,\ldots,n\}$, as claimed.

So we find that $q^{**}$ belongs to the set
\begin{multline*}
\mathcal{Q}^{**}:=\{q\in\Sigma_\mathcal{X}: q(x_1)=p(x_1)+\delta,~q(x_i)\leq p(x_i)\text{ for all }i\in\{2,\ldots,r\}\\\text{ and }q(x_i)=0\text{ for all }i\in\{r+1,\ldots,n\}\}.
\end{multline*}
Since $\min_{q\in\mathcal{Q}^*} E_q(f)=E_{q^{**}}(f)$ and $\mathcal{Q}^{**}\subseteq\mathcal{Q}^*$, this implies that
\begin{equation*}
\min_{q\in\mathcal{Q}^{**}} E_q(f)=\min_{q\in\mathcal{Q}^*}E_q(f)=\underline{E}^\mathrm{TV}_{p,\delta}(f).
\end{equation*}
Since the set of all $q\in\mathcal{Q}^{**}$ that attain the minimum $\min_{q\in\mathcal{Q}^{**}} E_q(f)$ is closed, we can choose $q_\mathrm{TV}$ to be one among them that minimizes $q_\mathrm{TV}(x_r)$.

Finally, we show that $q_\mathrm{TV}(x_i)=p(x_i)$ for all $i\in\{2,\ldots,r-1\}$. Assume \emph{ex absurdo} that there is some $i\in\{2,\ldots,r-1\}$ such that $q_\mathrm{TV}(x_i)<p(x_i)$. Consider any such index $i$. Since $q_\mathrm{TV}(x_i)\leq p(x_i)$ for all $i\in\{2,\ldots,r\}$, this implies that $\sum_{i=2}^{r-1} q_\mathrm{TV}(x_i)<\sum_{i=2}^{r-1} p(x_i)$. Since $q_\mathrm{TV}(x_1)=p(x_1)+\delta$ and $q_\mathrm{TV}(x_i)=0$ for all $i\in\{r+1,\ldots,n\}$, this implies that
\begin{align*}
q_{\mathrm{TV}}(x_r)
&=1-q_{\mathrm{TV}}(x_1)-\sum_{i=2}^{r-1} q_{\mathrm{TV}}(x_i)\\
&>1-p(x_1)-\delta-\sum_{i=2}^{r-1} p(x_i)\\
&\geq1-p(x_1)-\sum_{i=r+1}^n-\sum_{i=2}^{r-1}p(x_i)=p(x_r)\geq0.
\end{align*}
So $q_\mathrm{TV}(x_i)<p(x_i)$ but $q_{\mathrm{TV}}(x_r)>0$. Then we can construct a new probability mass function $q'$ by moving mass $\alpha=\min\{p(x_i)-q_\mathrm{TV}(x_i),q_\mathrm{TV}(x_r)\}>0$ from $x_r$ to $x_i$. This mass function $q'$ still belongs to $\mathcal{Q}^{**}$, and the expectation of $f$ does not increase since $f(x_i)\leq f(x_r)$. This however contradicts the definition of $q_\mathrm{TV}$, which was chosen to have the lowest value of $q(x_r)$ among all $q\in\mathcal{Q}^{**}$ that attain the minimum $\min_{q\in\mathcal{Q}^{**}} E_q(f)$. So we must have that $q_\mathrm{TV}(x_i)=p(x_i)$ for all $i\in\{2,\ldots,r-1\}$, as claimed.

So we conclude that our chosen $q_\mathrm{TV}\in\mathcal{Q}^{**}\subseteq\mathcal{Q}^*\subseteq B_\mathrm{TV}(p,\delta)$ such that $\underline{E}^\mathrm{TV}_{p,\delta}(f)=E_{q_\mathrm{TV}}(f)$ satisfies $q_\mathrm{TV}(x_1)=p(x_1)+\delta$, $q_\mathrm{TV}(x_i)=p(x_i)$ for all $i\in\{2,\ldots,r-1\}$ and $q_\mathrm{TV}(x_i)=0$ for all $i\in\{r+1,\ldots,n\}$, and therefore also
\begin{equation*}
q_\mathrm{TV}(x_r)=1-q_\mathrm{TV}(x_1)-\sum_{i=2}^{r-1} q_\mathrm{TV}(x_i)=1-p(x_1)-\delta-\sum_{i=2}^{r-1} p(x_i)=\sum_{i=r}^{n} p(x_i)-\delta>0,
\end{equation*}
as required.
\end{proof}


\end{document}
